% ============================================================
%  Section 1 --- Introduction
% ============================================================

\section{Introduction}
\label{sec:intro}

Les architectures microservices déployées sur Kubernetes sont devenues la norme dans les
organisations qui adoptent une cadence de livraison continue. Chaque service peut être mis
à jour, mis à l'échelle ou reconfiguré indépendamment, ce qui produit un flux permanent
de changements bénins dans le signal d'observabilité : déploiements progressifs, autoscaling
horizontal, modifications de quota. Les outils de détection d'anomalies classiques --- fondés
sur des seuils statiques ou des z-scores univariés \cite{myrtollari2025} --- ne distinguent pas
ces drifts bénins des vraies pannes. Il en résulte un niveau élevé de faux positifs en production, qui épuise la
confiance des équipes d'exploitation et retarde la réaction aux incidents réels.

EWAT (\emph{Early Warning and Anomaly Typing}) est un pipeline de détection précoce et de
typage automatique des anomalies conçu pour répondre à ce problème. Son positionnement est
explicitement distinct du RCA (\emph{Root Cause Analysis}) \cite{fu2025} : là où le RCA est post-mortem
(où, pourquoi, après la panne), EWAT est un système d'alerte précoce (quoi, dans combien de
temps, avant la panne). L'objectif n'est pas d'identifier la cause racine, mais de détecter
les signes précurseurs d'une dégradation imminente et d'en caractériser le type pour orienter
la réaction opérationnelle.

\subsection{Contexte}

Ce stage s'est déroulé chez Devoteam, société de conseil spécialisée en transformation
numérique, dans un contexte de montée en puissance de l'observabilité des plateformes cloud.
Le cluster expérimental \texttt{observit-cluster1} (RKE2 v1.32, 9 nœuds) héberge l'application
\emph{Online Boutique} de Google, une boutique en ligne de référence composée de microservices
hétérogènes. La stack d'observabilité combine Prometheus pour les métriques d'infrastructure
et un OpenTelemetry Collector pour les traces et logs applicatifs --- deux sources
complémentaires qui constituent ensemble le signal $S(t) \in \mathbb{R}^{N \times 17}$.

Les injections de pannes contrôlées sont réalisées avec Chaos Mesh, un opérateur Kubernetes
qui permet de simuler des scénarios reproductibles : saturation CPU, fuite mémoire,
déploiement défectueux, montée en charge soudaine. Cet outillage permet de constituer un
dataset labellisé tout en conservant un cluster réel, avec ses couplages inter-services et ses
bruits opérationnels --- une condition indispensable pour que les résultats soient transférables
en production.

\subsection{Problème}

Un déploiement progressif (\texttt{rolling\_deploy}) et une saturation CPU (\texttt{cpu\_starvation})
produisent tous deux des anomalies statistiques dans les métriques de latence et d'erreur.
Un z-score à $\sigma = 2.5$ déclenche une alerte dans les deux cas sans distinction.
Ce comportement est documenté dans la section~\ref{sec:alerts_results} : quel que soit
le seuil $\sigma \in [2.0, 3.5]$, le taux de fausses alarmes sur les drifts bénins reste
à 100\,\%.

La distinction fondamentale entre drift et anomalie exige une modélisation explicite de
quatre régimes opérationnels $\theta(t) \in \Theta$ plutôt que trois.
Trois régimes (\emph{normal}, \emph{drift}, \emph{anomalie}) sont insuffisants car ils
ne modélisent pas le cas hybride $\theta_{\text{drift} \cap \text{anomaly}}$ : un
déploiement défectueux qui présente simultanément les signatures d'un drift bénin
(changement de distribution graduel) et d'une anomalie réelle (dégradation des SLOs).
Ce quatrième régime, incarné par le scénario \texttt{faulty\_deploy\_overlap} du dataset,
requiert un mécanisme de \emph{look-through} pour ne pas supprimer le signal d'anomalie
pendant la transition de drift.

\subsection{Contributions}

Ce travail apporte quatre contributions principales :

\begin{enumerate}
  \item \textbf{Formalisation à quatre régimes et pipeline EWAT \textit{end-to-end}.}
    Définition formelle de $G(t)$, $S(t) \in \mathbb{R}^{N \times 17}$,
    des quatre régimes $\Theta$, et d'un pipeline en quatre étapes séquentielles
    (Section~\ref{sec:formalisation}). Implémentation complète avec 302 tests unitaires
    et lint propre (Ruff, mypy).

  \item \textbf{H1 validée --- structurabilité empirique des types de pannes.}
    Les embeddings STGCN forment des clusters bien séparés dans l'espace latent :
    silhouette test $= 0.519 \pm 0.092$ sur 5 graines, minimum $= 0.414 \gg 0.3$
    (seuil de Kaufman \& Rousseeuw \cite{kaufman1990}).
    $K = 10$ clusters stables, $K^* \in \{9, 10, 11\}$ selon la graine.

  \item \textbf{H3 validée --- précurseurs typés à horizon $k^*$ opérationnel.}
    AUROC moyen $= 0.973 \pm 0.012$ sur 5 graines, 7--8 types prédictibles sur $K$.
    Horizon optimal $k^* = 6$ steps (3\,min) pour 5 types sur 8 --- lead time opérationnel
    directement exploitable.

  \item \textbf{H2 négatif exploitable --- motivation de ewat\_v4.}
    Le mécanisme look-through n'apporte pas de réduction significative du FPR sur anomalie
    ($p = 0.27$). La cause est identifiée précisément : les épisodes de $\approx 21$ steps
    laissent $\approx 7$ steps utiles après le warm-up du DriftDetector (8 steps) et la
    confirmation temporelle (3--6 steps). Ce résultat négatif motive la collecte ewat\_v4
    avec des épisodes $\geq 40$ steps.
\end{enumerate}

\subsection{Plan du document}

La section~\ref{sec:formalisation} pose les bases mathématiques du problème.
La section~\ref{sec:dataset} décrit le dataset ewat\_v3 (299 épisodes, 15 scénarios).
La section~\ref{sec:pipeline} détaille l'implémentation de chaque étape du pipeline.
La section~\ref{sec:results} rapporte les résultats des hypothèses H1, H2a, H2b et H3,
ainsi que la comparaison avec la baseline z-score.
La section~\ref{sec:experiments} présente les expériences complémentaires : ablation des
modalités, comparaison d'architectures, ontologie causale et transfert de domaine sur RCAEval.
La section~\ref{sec:discussion} discute les apports, les limitations et les perspectives
pour ewat\_v4.
